%==================================================================
% PENGATURAN UTAMA DOKUMEN
% Template Tugas Akhir Politeknik Negeri Banyuwangi
% Acuan: Pedoman Mutu Tugas Akhir Poliwangi 2023, Bab 6
%
% FILE INI TIDAK PERLU DIUBAH.
% Kalau ingin mengubah format, ubah dengan sadar dan catat alasannya.
%==================================================================

%------------------------------------------------------------------
% 1. Bahasa dan pengkodean huruf
%------------------------------------------------------------------
\usepackage[T1]{fontenc}
\usepackage[utf8]{inputenc}
\usepackage[indonesian]{babel}

%------------------------------------------------------------------
% 2. Huruf Times New Roman, jarak baris 1,5
%    (Pedoman 6.1 butir 2)
%------------------------------------------------------------------
\usepackage{mathptmx}     % huruf Times + rumus matematika bergaya Times
\usepackage{setspace}
\setstretch{1.5}          % jarak baris isi = 1,5 spasi

%------------------------------------------------------------------
% 3. Kertas A4 dan margin
%    (Pedoman Lampiran 1: kiri 3 cm, atas 2 cm, kanan 2 cm, bawah 2 cm)
%------------------------------------------------------------------
\usepackage[a4paper,
            left=3cm,
            top=2cm,
            right=2cm,
            bottom=2cm,
            headsep=0.6cm,
            footskip=1.0cm]{geometry}

%------------------------------------------------------------------
% 4. Paragraf: baris pertama masuk 1,25 cm, jarak sebelum/sesudah 0
%    (Pedoman 6.1 butir 3)
%------------------------------------------------------------------
\usepackage{indentfirst}
\setlength{\parindent}{1.25cm}
\setlength{\parskip}{0pt}

% Mengurangi masalah perataan baris pada teks terjustifikasi
\tolerance=200
\emergencystretch=2.5em
\hyphenpenalty=200
\hbadness=3000

%------------------------------------------------------------------
% 5. Judul BAB, sub-bab, dan sub-sub-bab
%    (Pedoman Lampiran 2: "BAB 1" lalu judul, rata tengah, tebal)
%------------------------------------------------------------------
\usepackage{titlesec}

\titleformat{\chapter}[display]
  {\normalfont\fontsize{12}{14}\bfseries\centering}
  {BAB \thechapter}
  {0.6\baselineskip}
  {\MakeUppercase}
\titlespacing*{\chapter}{0pt}{-1.4cm}{1.2\baselineskip}

\titleformat{\section}
  {\normalfont\fontsize{12}{14}\bfseries}
  {\thesection}{0.6em}{}
\titlespacing*{\section}{0pt}{12pt}{6pt}

\titleformat{\subsection}
  {\normalfont\fontsize{12}{14}\bfseries}
  {\thesubsection}{0.6em}{}
\titlespacing*{\subsection}{0pt}{10pt}{4pt}

% Judul bab yang tidak bernomor (DAFTAR ISI, ABSTRAK, dst.)
\titleformat{name=\chapter,numberless}[block]
  {\normalfont\fontsize{12}{14}\bfseries\centering}
  {}{0pt}
  {\MakeUppercase}
\titlespacing*{name=\chapter,numberless}{0pt}{-1.4cm}{1.2\baselineskip}

%------------------------------------------------------------------
% 6. Keterangan tabel dan gambar
%    Tabel : nomor + judul di ATAS, rata kiri, tebal, spasi 1
%    Gambar: nomor + judul di BAWAH, rata tengah, tebal, spasi 1
%    (Pedoman 6.1 butir 4 dan 6)
%------------------------------------------------------------------
\usepackage[font={normalsize,stretch=1},
            labelfont=bf,
            textfont=bf,
            labelsep=space,
            singlelinecheck=false]{caption}
\usepackage{subcaption}

\captionsetup[table]{position=above, justification=raggedright, skip=6pt}
\captionsetup[figure]{position=below, justification=centering, skip=6pt}
\captionsetup[lstlisting]{position=above, justification=raggedright, skip=6pt}

% Penomoran tabel/gambar mengikuti nomor bab: Tabel 2.1, Gambar 3.2
\renewcommand{\thetable}{\arabic{chapter}.\arabic{table}}
\renewcommand{\thefigure}{\arabic{chapter}.\arabic{figure}}

%------------------------------------------------------------------
% 7. Rumus matematika
%    Nomor rumus mengikuti nomor bab dan rata kanan (Pedoman 6.1 butir 5)
%    Rumus dimulai dari batas kiri margin (opsi fleqn di laporan.tex)
%------------------------------------------------------------------
\usepackage{amsmath}
\numberwithin{equation}{chapter}
\setlength{\mathindent}{0pt}

%------------------------------------------------------------------
% 8. Daftar isi
%------------------------------------------------------------------
\usepackage{tocloft}

\renewcommand{\cftchapfont}{\normalfont\bfseries}
\renewcommand{\cftchappagefont}{\normalfont\bfseries}
\renewcommand{\cftchappresnum}{BAB }
\renewcommand{\cftchapaftersnum}{ }
\newlength{\lebarbab}
\settowidth{\lebarbab}{\bfseries BAB 1\ }
\cftsetindents{chap}{0pt}{\lebarbab}

\setlength{\cftbeforechapskip}{0pt}
\setlength{\cftbeforesecskip}{0pt}
\setlength{\cftbeforesubsecskip}{0pt}
\renewcommand{\cftsecleader}{\cftdotfill{\cftdotsep}}
\setlength{\cftsecindent}{0pt}
\setlength{\cftsecnumwidth}{2.6em}
\setlength{\cftsubsecindent}{2.6em}
\setlength{\cftsubsecnumwidth}{3.2em}
\setlength{\cftfigindent}{0pt}
\setlength{\cfttabindent}{0pt}
\renewcommand{\cftdotsep}{1}

% Daftar isi, daftar tabel, dan daftar gambar memakai spasi 1
\renewcommand{\cfttoctitlefont}{\hfill\normalfont\bfseries}
\renewcommand{\cftaftertoctitle}{\hfill}
\renewcommand{\cftloftitlefont}{\hfill\normalfont\bfseries}
\renewcommand{\cftafterloftitle}{\hfill}
\renewcommand{\cftlottitlefont}{\hfill\normalfont\bfseries}
\renewcommand{\cftafterlottitle}{\hfill}

%------------------------------------------------------------------
% 9. Tabel dan gambar
%------------------------------------------------------------------
\usepackage{graphicx}
\graphicspath{{gambar/}}
\usepackage{float}
\usepackage{tabularx}
\usepackage{multirow}
\usepackage{longtable}
\usepackage{rotating}       % halaman landscape
\usepackage{array}

%------------------------------------------------------------------
% 10. Daftar butir
%------------------------------------------------------------------
\usepackage{enumitem}
\setlist{nosep}

\newenvironment{packed_enum}
  {\begin{enumerate}[leftmargin=1.5\parindent]}
  {\end{enumerate}}
\newenvironment{packed_item}
  {\begin{itemize}[leftmargin=1.5\parindent]}
  {\end{itemize}}

%------------------------------------------------------------------
% 11. Daftar pustaka gaya APA (penulis, tahun)
%------------------------------------------------------------------
\usepackage[round,authoryear]{natbib}
\bibliographystyle{apalike}

%------------------------------------------------------------------
% 12. Kode program
%------------------------------------------------------------------
\usepackage{listings}
\usepackage[table]{xcolor}   % opsi table: untuk \rowcolor dan \cellcolor

\definecolor{codegreen}{rgb}{0,0.5,0}
\definecolor{codegray}{rgb}{0.45,0.45,0.45}
\definecolor{codepurple}{rgb}{0.55,0,0.65}
\definecolor{backcolour}{rgb}{0.97,0.97,0.95}

\lstdefinestyle{poliwangi}{
    backgroundcolor=\color{backcolour},
    commentstyle=\color{codegreen},
    keywordstyle=\color{magenta}\bfseries,
    numberstyle=\tiny\color{codegray},
    stringstyle=\color{codepurple},
    basicstyle=\ttfamily\fontsize{9}{10.8}\selectfont,
    columns=fullflexible,
    breakatwhitespace=false,
    breaklines=true,
    captionpos=t,
    keepspaces=true,
    numbers=left,
    numbersep=6pt,
    showspaces=false,
    showstringspaces=false,
    showtabs=false,
    tabsize=2,
    frame=single,
    rulecolor=\color{codegray},
    xleftmargin=1.2em,
    framexleftmargin=1.2em
}
\lstset{style=poliwangi}

\AtBeginDocument{%
  \renewcommand{\lstlistingname}{Kode}%
  \renewcommand{\thelstlisting}{\arabic{chapter}.\arabic{lstlisting}}%
}

%------------------------------------------------------------------
% 13. Gambar dengan TikZ (diagram alir, grafik)
%------------------------------------------------------------------
\usepackage{tikz}
\usetikzlibrary{shapes.geometric, arrows.meta, positioning, calc}
\usepackage{pgfplots}
\pgfplotsset{compat=1.18}

\tikzset{
  kotak/.style   = {rectangle, rounded corners, minimum width=3cm,
                    minimum height=0.9cm, text centered, draw=black,
                    align=center, font=\small},
  proses/.style  = {kotak, fill=blue!8},
  mulai/.style   = {kotak, fill=green!12},
  keputusan/.style = {diamond, aspect=2.2, draw=black, align=center,
                      font=\small, fill=orange!12, inner sep=1pt},
  panah/.style   = {thick, -{Stealth[length=2.5mm]}}
}

%------------------------------------------------------------------
% 14. Tautan dan rujukan otomatis (Gambar 2.1, Tabel 3.2, persamaan 4.1)
%------------------------------------------------------------------
\usepackage[hidelinks]{hyperref}
\usepackage{cleveref}

\crefname{figure}{Gambar}{Gambar}
\Crefname{figure}{Gambar}{Gambar}
\crefname{table}{Tabel}{Tabel}
\Crefname{table}{Tabel}{Tabel}
\crefname{equation}{persamaan}{persamaan}
\Crefname{equation}{Persamaan}{Persamaan}
\crefname{lstlisting}{Kode}{Kode}
\Crefname{lstlisting}{Kode}{Kode}

%------------------------------------------------------------------
% 15. Perintah bantu
%------------------------------------------------------------------
\usepackage{ifthen}

% Halaman kosong diberi keterangan, bukan halaman benar-benar kosong
\newcommand{\halamankosong}{%
  \clearpage
  \thispagestyle{empty}%
  \vspace*{\fill}%
  \begin{center}\textit{``Halaman ini sengaja dikosongkan''}\end{center}%
  \vspace*{\fill}%
  \clearpage
}

% Jarak bantu: \jarak{2} = 2 x 1 spasi
\newcommand{\jarak}[1]{\vspace{#1\baselineskip}}

% Bagian awal (ABSTRAK, KATA PENGANTAR, dst.) dimasukkan ke daftar isi
% TANPA awalan "BAB".
\makeatletter
\newcommand{\l@front}[2]{%
  \par\noindent\cftchapfont #1\nobreak
  \leaders\hbox{$\m@th\mkern\cftdotsep mu\hbox{.}\mkern\cftdotsep mu$}%
  \hfill\nobreak{\cftchappagefont #2}\par}
\makeatother

\newcommand{\frontmatterentry}[1]{%
  \phantomsection
  \addcontentsline{toc}{front}{#1}%
}

% Lampiran: bernomor urut sendiri dan otomatis masuk DAFTAR LAMPIRAN
\newcounter{lampiran}
\makeatletter
\newcommand{\l@lampiran}[2]{%
  \par\noindent LAMPIRAN \hb@xt@0.6em{}#1\nobreak
  \leaders\hbox{$\m@th\mkern\cftdotsep mu\hbox{.}\mkern\cftdotsep mu$}%
  \hfill\nobreak #2\par}
\makeatother

\newcommand{\lampiran}[1]{%
  \clearpage
  \refstepcounter{lampiran}%
  \addcontentsline{loa}{lampiran}{\thelampiran\hspace{0.8em}#1}%
  \phantomsection
  \begin{center}
    \textbf{LAMPIRAN \thelampiran}\par
    \vspace{0.2\baselineskip}
    \textbf{#1}\par
  \end{center}
  \vspace{0.4\baselineskip}%
}

\makeatletter
\newcommand{\daftarlampiran}{%
  \chapter*{DAFTAR LAMPIRAN}%
  \frontmatterentry{DAFTAR LAMPIRAN}%
  \begin{spacing}{1}%
  \@starttoc{loa}%
  \end{spacing}%
}
\makeatother

%==================================================================
% AKHIR PENGATURAN
%==================================================================
